% Pursuing Stacks (cotes 134-2 to 134-8): archive cote 134-3, site page 91 (PDF page 92).
% notes for ch. II. Draft compared: ps2.tex:595-606 ; ps2.tex:708-724 ; ps2.tex:761-773 ; ps2.tex:960-974.
% AI-transcribed by Claude Opus 5.5 (claude-opus-5-5, adaptive thinking, effort high), Message Batches API, batch ps1 (msgbatch_01Rv5hCaNTtUt4JHx1yqoHPL), 2026-11-12; not yet checked by eye.
\scanpage{134-3/91}
II iii

($^{13}$) I am relying here on Quillen's set-up as the \deleted{new} only axiomatic
set-up for axiomatic homotopy theory which was known to me at th\over{\unsure{a}}{a}t moment\unsure{.}
In the cour\over{\unsure{s}}{s}e of my re\over{\unsure{fl}}{fl}ection, \deleted{a} different set-ups \over{\unsure{a}}{a}re peeling out,
which seem more directly suited for finding "modelizers" for the specific
category (Hot) (notably, the so-called \uline{asphericity structures} and, still
more importantly, the \uline{contractibility structures}). Thus, the notion
of a model category in Quillen's sense is somewhat fading into the back-%
ground after section 2\over{\unsure{5}}{3}, where one \over{\unsure{m}}{m}ain new idea \xout{appears}\inserted[typed]{gets emphasized}\hand{,} namely the
description of a canonical functor
\[ M \hand{\longrightarrow} (\mathrm{Hot}) \]
for any category M endowed with a structure of a (locally contractible)
"site"...

($^{14}$) Cf. Jean Giraud, Méthode de la Descente, Bull. Soc. Math.
Mémoire n\textsuperscript{o}2 (dec. 1964).

($^{15}$) It turns out that the \inserted[typed]{relevant\hand{,}} topology is \xout{rarely} \inserted[typed]{not always} the canonical one -
however the intuition which appears here is going to be basic in the
whole "Modelizing Story"; \deleted{\unsure{a}} compare note ($^{13}$) above.

($^{16}$) It will turn out later that these kinds of complexes, as we\over{\unsure{l}}{l}l
as the multicomplexes envisioned below (eve\over{\unsure{n}}{n} the "mixed" ones, namely
semisim\deleted{\unsure{m}}plicial or simplicial for some indices, and cubical (\over{\unsure{W}}{w}ith or
without symmetries)\hand{,} or even hemispherical\hand{,} \xout{in} \inserted[typed]{for} some others,..\over{\unsure{a}}{)} do indeed
model homotopy types. This is a consequence of the criterion for test
categories (section 31), and of the fact that a product of test catego-%
ries is again a test category (cf. section 74).

($^{1\over{\unsure{6}}{7}}$) The notes are a little floppy here - the "cells" are going to be
compact \over{\unsure{o}}{o}nly under suitable restriction on the ordered set K, for instan-%
ce if K is finite - however, it is clear I want also to consider arbitra-%
ry ord\over{\unsure{rr}}{er}ed sets.

($^{1\over{\unsure{7}}{8}}$) It does make a sense however to ask \over{\unsure{wh}}{wh}ether this categor\over{\unsure{y}}{y} is a
"modelizer" (in the preci\over{\unsure{s}}{z}e sense of section 28) - I believe it is, but
did'nt check it yet.

($^{1\over{\unsure{8}}{9}}$) As we'll see, in the all-important case $\underline{B}$ = (Cat), the canonical
topolog\over{\unsure{y}}{y} is definitely \uline{not} the right one !

($^{\over{\unsure{19}}{20}}$) The consideration of homotopy types "of finite presentation" turns
out to be irrelevant at this stage o\over{\unsure{n}}{f} the reflection, leading up mainly
to the notion of "test categories" and "test functors", which are inde-%
pendant of "finite type" kind of conditions.

($^{2\over{\unsure{0}}{1}}$) For the time being, these considerations of cellular structures

\mspage{91}
